\subsection{模关于理想的重数}

设 M 是有限生成的 A-模，q 是 A 的一个理想，包含于 A 的根基中，且 M/qM 具有有限长度。假设 M 不为 0，置 \(d = \dim_{\mathbf{A}}(\mathbf{M})\)。根据第 § 4 节第 3 号定理 2 的推论及第 4 号注 2，存在唯一的整数 \(e_{\mathfrak{q}}(\mathbf{M}) > 0\)，使得对于每个整数 \(n \geqslant 1\)

\[
\operatorname{long} _ {\mathbf {A}} (\mathrm{M} / \mathfrak {q} ^ {n + 1} \mathrm{M}) = e _ {\mathfrak {q}} (\mathrm{M}) \frac {n ^ {d}}{d !} + \beta_ {n} n ^ {d - 1}
\]

其中 \(\beta_{n}\) 在 \(n\) 无限增大时趋于一个极限。

定义 1. --- 整数 \(e_{\mathfrak{q}}(\mathbf{M})\) 称为 A-模 \(\mathbf{M}\) 关于理想 q 的重数。

当需要指明环 A 时，也记作 \(e_{\mathfrak{q}}^{\mathbf{A}}(\mathbf{M})\)。当 A 是极大理想为 m 的局部环时，分别用 \(e(\mathbf{M})\) 或 \(e^{\mathbf{A}}(\mathbf{M})\) 表示 \(e_{\mathfrak{m}}(\mathbf{M})\) 或 \(e_{\mathfrak{m}}^{\mathbf{A}}(\mathbf{M})\)。

注记。--- 1) 若 q' 是 A 的一个理想，包含于 A 的根基中且包含 q，则 \(e_{\mathfrak{q}'}(\mathbf{M}) \leqslant e_{\mathfrak{q}}(\mathbf{M})\)；并且，若 M 的 \(q'\)-进滤过是 q-良好的，则有 \(e_{\mathfrak{q}'}(\mathbf{M}) = e_{\mathfrak{q}}(\mathbf{M})\)（§ 4，第 3 号，定理 2）。

\begin{enumerate}
\def\labelenumi{\arabic{enumi})}
\setcounter{enumi}{1}
\item
  若 M 具有有限长度，则有 \(e_{\mathfrak{q}}(\mathbf{M}) = \mathrm{long}_{\mathbf{A}}(\mathbf{M})\)（§ 4，第 3 号，注 3）。
\item
  若 \(d > 0\)，则有
\end{enumerate}

\[
\operatorname{long} _ {\mathrm{A} / \mathfrak {q}} (\mathfrak {q} ^ {n} \mathrm{M} / \mathfrak {q} ^ {n + 1} \mathrm{M}) = e _ {\mathfrak {q}} (\mathrm{M}) \frac {n ^ {d - 1}}{(d - 1) !} + \alpha_ {n} n ^ {d - 2}
\]

其中 \(\alpha_{n}\) 在 \(n\) 无限增大时趋于一个极限（§ 4，第 3 号，定理 2 的推论）。

\begin{enumerate}
\def\labelenumi{\arabic{enumi})}
\setcounter{enumi}{3}
\tightlist
\item
  由于根据 § 4，第 4 号定理 3 的推论，有
\end{enumerate}

\[
e _ {\mathrm{q}} (\mathbf {M}) = \sum e _ {\mathrm{q} _ {\mathrm{m}}} (\mathbf {M} _ {\mathrm{m}})
\]

求和遍历 A 的所有极大理想 m，其中满足

\[
\mathfrak {m} \in \operatorname{Supp} (\mathbf {M}) \cap \mathrm{V} (\mathfrak {q}) \quad \text { and } \quad \dim_ {\mathrm{A} _ {\mathfrak {m}}} (\mathrm{M} _ {\mathfrak {m}}) = d  .
\]

由注记 2 和 3 可知，\(e_{\mathfrak{q}}(\mathbf{M})\) 只取决于分次 A/q-模 \(\mathrm{gr}_{\mathfrak{q}}(\mathbf{M})\)。因此：

命题 1.--- 设 \(\hat{\mathbf{A}}\) 和 \(\hat{\mathbf{M}}\) 分别是 \(\mathbf{A}\) 和 \(\mathbf{M}\) 关于 q-进拓扑的完备化；则 \(e_{\mathfrak{q}}^{\mathbf{A}}(\mathbf{M}) = e_{\mathfrak{q}\hat{\mathbf{A}}}^{\hat{\mathbf{A}}}(\hat{\mathbf{M}})\)。

命题 2. --- 假设 A 是正则局部环（§ 5，第 1 号，定义 1）；则 e(A) = 1。

这由 § 5 第 2 号的定理 1 得出。

注记 5. --- 即使 A 不正则，也可能有 \(e(A) = 1\)（p.~104，练习 5）。事实上，诺特局部环 A 正则，当且仅当 \(\hat{A}\) 是整环且有 \(e(A) = 1\)（p.~108，练习 24）。

例。--- 根据定义，有 \(e_{\mathfrak{q}^{r}}(\mathbf{M}) = r^{d}e_{\mathfrak{q}}(\mathbf{M})\)，其中 \(d = \dim_{\mathbf{A}}(\mathbf{M})\)。因此，若 A 是正则局部环，则 \(e_{\mathfrak{m}_{\mathbf{A}}^{r}}(\mathbf{A}) = r^{d}\)。例如，若 A 是离散赋值环，则 \(e_{\mathfrak{q}}(\mathbf{A}) = \text{long}(\mathbf{A}/\mathfrak{q})\)。

设 q 是 A 的一个理想，包含于 A 的根基中，令 \(\mathcal{C}(\mathfrak{q})\) 为满足 \(M/qM\) 具有有限长度的有限生成 A-模的同构类集合。对于每个 \(d \in \mathbf{N}\)，以 \(\mathcal{C}(\mathfrak{q})_{\leq d}\) 表示 \(\mathcal{C}(\mathfrak{q})\) 中由维数 \(\leqslant d\) 的 A-模的同构类组成的部分。定义映射 \(e_{\mathfrak{q},d}: \mathcal{C}(\mathfrak{q})_{\leq d} \to \mathbf{Z}\)：规定 \(e_{\mathfrak{q},d}(\mathbf{M}) = e_{\mathfrak{q}}(\mathbf{M})\)（当 \(\dim(\mathbf{M}) = d\) 时），否则 \(e_{\mathfrak{q},d}(\mathbf{M}) = 0\)。根据第 § 4 节第 3 号的命题 5，这个映射是可加的；由此（A, VIII, § 10, 第 2 号）得到一个仍记作 \(e_{\mathfrak{q},d}\) 的同态，它从 Grothendieck 群 \(K(\mathcal{C}(\mathfrak{q})_{\leq d})\) 到 \(\mathbf{Z}\)，并且在 \(K(\mathcal{C}(\mathfrak{q})_{\leq d-1})\) 上为零。仿照 § 1 第 5 号的论证，得到：

命题 3. — 设 \(\mathbf{M}\) 是有限生成的 A-模，维数 \(d \geqslant 0\)。令 \(\Phi\) 为由极小元 \(\mathfrak{p}\) 构成的集合，这些极小元属于 \(\operatorname{Supp}(\mathbf{M})\) 且满足 \(\dim(\mathbf{A}/\mathfrak{p}) = d\)。设 \(\mathfrak{q}\) 是 A 的一个理想，包含于 A 的根基中，且 \(\mathbf{M}/\mathfrak{q}\mathbf{M}\) 具有有限长度。则有

\[
e _ {\mathfrak {q}} (\mathrm{M}) = \sum_ {\mathfrak {p} \in \Phi} \operatorname{long} _ {\mathrm{A} _ {\mathfrak {p}}} (\mathrm{M} _ {\mathfrak {p}}) \cdot e _ {\mathfrak {q}} (\mathrm{A} / \mathfrak {p}).
\]

推论。— 假设 A 是半局部环，且 q 是 A 的一个定义理想。

\begin{enumerate}
\def\labelenumi{\alph{enumi})}
\item
  有 \(e_{\mathfrak{q}}(\mathbf{A}) = \sum_{\mathfrak{p}} e_{\mathfrak{q}}(\mathbf{A}/\mathfrak{p})\)，其中 \(\mathfrak{p}\) 遍历 \(\mathbf{A}\) 的极小素理想中满足 \(\dim(\mathbf{A}/\mathfrak{p}) = \dim(\mathbf{A})\) 的那些理想。
\item
  假设 A 是整环，且 M 是有限生成的 A-模，满足 \(\dim_{\mathbf{A}}(\mathbf{M}) = \dim(\mathbf{A})\)。则有 \(e_{\mathfrak{q}}(\mathbf{M}) = \operatorname{rg}(\mathbf{M}) \cdot e_{\mathfrak{q}}(\mathbf{A})\)。
\end{enumerate}

